\documentclass[11pt,fontset=fandol]{ctexart}
\usepackage[a4paper,margin=23mm]{geometry}
\usepackage{amsmath,amssymb,mathtools,bm,mathrsfs,esint}
\usepackage{graphicx,longtable,booktabs,array,calc}
\usepackage[normalem]{ulem}
\usepackage{xcolor,hyperref}
\hypersetup{hidelinks}
\setcounter{secnumdepth}{-1}
\newcounter{none}
\providecommand{\tightlist}{\setlength{\itemsep}{0pt}\setlength{\parskip}{0pt}}
\providecommand{\passthrough}[1]{#1}
\usepackage{amssymb,mathrsfs,esint,bm,newunicodechar}
\allowdisplaybreaks
\setlength{\emergencystretch}{3em}
\providecommand{\llangle}{\langle\!\langle}
\providecommand{\rrangle}{\rangle\!\rangle}
\newunicodechar{∞}{\ensuremath{\infty}}
\newunicodechar{∈}{\ensuremath{\in}}
\newunicodechar{∑}{\ensuremath{\sum}}
\newunicodechar{∫}{\ensuremath{\int}}
\newunicodechar{≤}{\ensuremath{\leq}}
\newunicodechar{≥}{\ensuremath{\geq}}
\newunicodechar{≠}{\ensuremath{\neq}}
\newunicodechar{→}{\ensuremath{\to}}
\newunicodechar{⇒}{\ensuremath{\Rightarrow}}
\newunicodechar{α}{\ensuremath{\alpha}}
\newunicodechar{β}{\ensuremath{\beta}}
\newunicodechar{γ}{\ensuremath{\gamma}}
\newunicodechar{λ}{\ensuremath{\lambda}}
\newunicodechar{φ}{\ensuremath{\varphi}}
\newunicodechar{ε}{\ensuremath{\varepsilon}}
\newunicodechar{σ}{\ensuremath{\sigma}}
\newunicodechar{θ}{\ensuremath{\theta}}
\newunicodechar{δ}{\ensuremath{\delta}}
\newunicodechar{η}{\ensuremath{\eta}}
\newunicodechar{π}{\ensuremath{\pi}}
\newunicodechar{ξ}{\ensuremath{\xi}}
\newunicodechar{Σ}{\ensuremath{\Sigma}}
\newunicodechar{Δ}{\ensuremath{\Delta}}
\newunicodechar{∂}{\ensuremath{\partial}}
\newunicodechar{∀}{\ensuremath{\forall}}
\newunicodechar{∃}{\ensuremath{\exists}}
\newunicodechar{√}{\ensuremath{\surd}}
\newunicodechar{−}{\ensuremath{-}}
\newunicodechar{μ}{\ensuremath{\mu}}
\newunicodechar{ρ}{\ensuremath{\rho}}
\newunicodechar{τ}{\ensuremath{\tau}}
\newunicodechar{ω}{\ensuremath{\omega}}
\newunicodechar{Ω}{\ensuremath{\Omega}}
\newunicodechar{χ}{\ensuremath{\chi}}
\newunicodechar{ψ}{\ensuremath{\psi}}
\newunicodechar{Γ}{\ensuremath{\Gamma}}
\newunicodechar{Φ}{\ensuremath{\Phi}}
\newunicodechar{Λ}{\ensuremath{\Lambda}}
\newunicodechar{ℝ}{\ensuremath{\mathbb{R}}}
\newunicodechar{ℤ}{\ensuremath{\mathbb{Z}}}
\newunicodechar{ℕ}{\ensuremath{\mathbb{N}}}
\newunicodechar{△}{\ensuremath{\triangle}}
\newunicodechar{⊗}{\ensuremath{\otimes}}
\begin{document}
\section{2023年杭州电子科技大学601数学分析考研真题解答(手写详细版)}\label{ux5e74ux676dux5ddeux7535ux5b50ux79d1ux6280ux5927ux5b66601ux6570ux5b66ux5206ux6790ux8003ux7814ux771fux9898ux89e3ux7b54ux624bux5199ux8be6ux7ec6ux7248}

来源：\texttt{HDU/601数学分析/2023/2023年杭州电子科技大学601数学分析考研真题解答(手写详细版).pdf}。原
PDF 页序：1, 2, 3, 4, 5, 6。

\section{2023年杭州电子科技大学数学分析手写详细版：逐页推导整理}\label{ux5e74ux676dux5ddeux7535ux5b50ux79d1ux6280ux5927ux5b66ux6570ux5b66ux5206ux6790ux624bux5199ux8be6ux7ec6ux7248ux9010ux9875ux63a8ux5bfcux6574ux7406}

\begin{quote}
依据原PDF第1---6页逐页本地OCR整理，保留手写版可核定的解题步骤。本稿由助手整理，不作为原作者逐字转录。标注``【OCR核正】''的部分由题面、上下文或计算恢复；原推导有缺口时，另标``助手补充修正''。
\end{quote}

原件：\href{/Users/oyamahappa/Documents/Agent/Codex/HDU/601数学分析/2023/2023年杭州电子科技大学601数学分析考研真题解答(手写详细版).pdf}{手写PDF}。互证：\href{/Users/oyamahappa/Documents/Agent/Codex/HDU/601数学分析/2023/2023年杭州电子科技大学601数学分析考研真题参考答案.pdf}{同年打印参考答案PDF}。

\subsection{原第1页：填空题1---5}\label{ux539fux7b2c1ux9875ux586bux7a7aux989815}

\subsubsection{一1.
迭代函数的导数}\label{ux4e001.-ux8fedux4ee3ux51fdux6570ux7684ux5bfcux6570}

本页明确询问 \(f'_{2023}(0)\)，其中 \(f(x)=x/\sqrt{1+x^2}\)，\(f_n\) 为
\(f\) 的 \(n\) 次复合。打印题面及打印答案当前对应函数值
\(f_{2023}(0)=0\)；两版差异分别保留。

由 \[
f_1(x)=\frac{x}{\sqrt{1+x^2}},\qquad
f_2(x)=\frac{x/\sqrt{1+x^2}}{\sqrt{1+x^2/(1+x^2)}}=\frac{x}{\sqrt{1+2x^2}},
\] 归纳得到 \[
f_n(x)=\frac{x}{\sqrt{1+nx^2}},\qquad
f'_{2023}(x)=\frac1{(1+2023x^2)^{3/2}}.
\] 因此 \(f'_{2023}(0)=1\)。【OCR核正】末式误识为 \(\theta\)
的自变量恢复为题目指定的零。

\subsubsection{一2. 根式极限}\label{ux4e002.-ux6839ux5f0fux6781ux9650}

保留有理化步骤： \[
\begin{aligned}
\lim_{n\to\infty}n(\sqrt{n^2+1}-n)
&=\lim_{n\to\infty}\frac{n[(n^2+1)-n^2]}{\sqrt{n^2+1}+n}\\
&=\lim_{n\to\infty}\frac1{\sqrt{1+1/n^2}+1}=\frac12.
\end{aligned}
\] 【OCR核正】有理化后的分子为 \(n\)，不能采用误识的 \(\pi\)。

\subsubsection{一3.
变上限积分的导数}\label{ux4e003.-ux53d8ux4e0aux9650ux79efux5206ux7684ux5bfcux6570}

设 \(f\) 在 \([a,b]\) 连续，\(F(x)=\int_a^x f(t)(t-x)\,dt\)。先分解： \[
F(x)=\int_a^x tf(t)\,dt-x\int_a^xf(t)\,dt.
\] 再求导： \[
F'(x)=xf(x)-\int_a^xf(t)\,dt-xf(x)=-\int_a^xf(t)\,dt.
\]

\subsubsection{一4. Gaussian 积分}\label{ux4e004.-gaussian-ux79efux5206}

原稿使用 Euler--Poisson 积分： \[
\int_0^\infty e^{-x^2}\,dx=\frac{\sqrt\pi}{2}.
\] 由偶性， \[
\int_{-\infty}^\infty e^{-x^2}\,dx
=2\int_0^\infty e^{-x^2}\,dx=\sqrt\pi.
\] 【OCR核正】本地OCR漏掉平方根，三处 \(\pi\)
按标准积分及同年打印答案恢复为 \(\sqrt\pi\)。

\subsubsection{一5.
分段函数的间断点}\label{ux4e005.-ux5206ux6bb5ux51fdux6570ux7684ux95f4ux65adux70b9}

\[
f(x)=\begin{cases}\cos(\pi x/2),&|x|\le1,\\|x-1|,&|x|>1.\end{cases}
\] 在 \(x=1\)，左右极限和函数值均为零，故连续。在 \(x=-1\)， \[
f((-1)^+)=0,\qquad f((-1)^-)=2,
\] 故唯一间断点为 \(-1\)，是跳跃间断点。其余各分支内部连续。

\subsection{原第2页：填空题6---8，计算题1---2}\label{ux539fux7b2c2ux9875ux586bux7a7aux989868ux8ba1ux7b97ux989812}

\subsubsection{一6. 收敛半径}\label{ux4e006.-ux6536ux655bux534aux5f84}

记 \(u_n=(1+1/n)^{n(n+1)}\)，则 \[
\rho=\lim_{n\to\infty}\sqrt[n]{|u_n|}
=\lim_{n\to\infty}(1+1/n)^{n+1}=e.
\] 所以 \(\sum_{n=2}^\infty u_nx^n\) 的收敛半径为 \(R=1/e\)。

\subsubsection{一7.
四分之一圆弧上的积分}\label{ux4e007.-ux56dbux5206ux4e4bux4e00ux5706ux5f27ux4e0aux7684ux79efux5206}

对第一象限圆弧 \(x^2+y^2=r^2\)，按半径 \(r>0\)，取
\(x=r\cos\theta,y=r\sin\theta\)、\(0\le\theta\le\pi/2\)。于是 \[
\int_Lxy\,ds=r^3\int_0^{\pi/2}\sin\theta\cos\theta\,d\theta
=\left[\frac{r^3}{2}\sin^2\theta\right]_0^{\pi/2}=\frac{r^3}{2}.
\]

\subsubsection{一8. 不定积分}\label{ux4e008.-ux4e0dux5b9aux79efux5206}

令 \(t=x^n\)，在 \(n\ne0\) 且被积函数有定义的区间内， \[
\begin{aligned}
\int\frac{dx}{x(1+x^n)}
&=\int\frac{x^{n-1}}{x^n(1+x^n)}\,dx
=\frac1n\int\frac{dt}{t(1+t)}\\
&=\frac1n\ln\left|\frac{t}{1+t}\right|+C
=\frac1n\ln\left|\frac{x^n}{1+x^n}\right|+C.
\end{aligned}
\]

\subsubsection{二1. 用 Green
公式计算上半圆积分}\label{ux4e8c1.-ux7528-green-ux516cux5f0fux8ba1ux7b97ux4e0aux534aux5706ux79efux5206}

计算沿上半圆 \(L\) 从 \((1,0)\) 到 \((-1,0)\) 的积分 \[
I=\int_L(-2xe^{-x^2}\sin y)\,dx+(e^{-x^2}\cos y+x^6)\,dy.
\] 令 \(L'\) 从 \((-1,0)\) 沿直线到 \((1,0)\)，\(D\)
为上半单位圆盘。原稿记闭曲线积分为 \(I_1\)、直线段积分为 \(I_2\)。由 \[
Q_x-P_y=(-2xe^{-x^2}\cos y+6x^5)-(-2xe^{-x^2}\cos y)=6x^5,
\] 有 \(I_1=6\iint_Dx^5\,dx\,dy=0\)。原稿的两种核算方式保留为 \[
6\int_0^\pi\cos^5\theta\,d\theta\int_0^1r^6\,dr=0,\qquad
6\int_{-1}^1x^5\sqrt{1-x^2}\,dx=0.
\] 在 \(L'\) 上 \(y=0,dy=0\)，故 \(I_2=0\)。因此 \(I=I_1-I_2=0\)。
【OCR核正】恢复极坐标 Jacobian、\(\cos^5\theta\)、圆盘上界
\(\sqrt{1-x^2}\)；按上述方向定义，组合关系应为减号。

\subsubsection{二2. 对数极限}\label{ux4e8c2.-ux5bf9ux6570ux6781ux9650}

由 \[
\ln(1+1/x)=\frac1x-\frac1{2x^2}+o(x^{-2}),\qquad x\to+\infty,
\] 得 \[
x-x^2\ln(1+1/x)=x-x+\frac12+o(1)\to\frac12.
\]

\subsection{原第3页：计算题3---5}\label{ux539fux7b2c3ux9875ux8ba1ux7b97ux989835}

\subsubsection{二3.
先积分辅助级数，再求导}\label{ux4e8c3.-ux5148ux79efux5206ux8f85ux52a9ux7ea7ux6570ux518dux6c42ux5bfc}

保留手写稿独有的辅助函数过程： \[
S(x)=\sum_{n=1}^\infty\frac{2n}{3^n}x^n,\qquad
h(x)=\sum_{n=1}^\infty\frac{2n}{3^n}x^{n-1},\qquad S(x)=xh(x).
\] 根值判别给出 \(\lim\sqrt[n]{2n/3^n}=1/3\)，故收敛半径为
\(3\)，两个端点通项不趋于零，收敛域为 \((-3,3)\)。在该区间逐项积分： \[
\int_0^x h(t)\,dt
=\sum_{n=1}^\infty\frac{2x^n}{3^n}=\frac{2x}{3-x}.
\] 于是 \[
S(x)=x\left(\frac{2x}{3-x}\right)'=\frac{6x}{(3-x)^2},
\qquad \sum_{n=1}^\infty\frac{2n}{3^n}=S(1)=\frac32.
\] 【OCR核正】根值判别的根号次数恢复为 \(n\)。

\subsubsection{二4.
直接参数化空间闭曲线}\label{ux4e8c4.-ux76f4ux63a5ux53c2ux6570ux5316ux7a7aux95f4ux95edux66f2ux7ebf}

\(\Gamma\) 为 \(x^2+y^2=1\) 与 \(x+z=1\) 的交线，从 \(z\)
轴正向看逆时针。令 \[
x=\cos\theta,\quad y=\sin\theta,\quad z=1-\cos\theta,\quad0\le\theta\le2\pi.
\] 则 \[
\begin{aligned}
\oint_\Gamma x^2y\,dx+y^2\,dy+z\,dz
&=\int_0^{2\pi}[-\cos^2\theta\sin^2\theta
+\sin^2\theta\cos\theta+(1-\cos\theta)\sin\theta]\,d\theta\\
&=-\frac18\int_0^{2\pi}(1-\cos4\theta)\,d\theta=-\frac\pi4.
\end{aligned}
\] 第二、三项各自整周期积分均为零。

\subsubsection{二5.
先投影，再用极坐标}\label{ux4e8c5.-ux5148ux6295ux5f71ux518dux7528ux6781ux5750ux6807}

\(V\) 由 \(z=\sqrt{x^2+y^2}\) 和 \(z=x^2+y^2\) 围成，其投影为
\(D:x^2+y^2\le1\)。依手写稿先对 \(z\) 积分： \[
\begin{aligned}
I&=\iiint_Vx^2\sqrt{x^2+y^2}\,dx\,dy\,dz\\
&=\iint_Dx^2\sqrt{x^2+y^2}
\left[\int_{x^2+y^2}^{\sqrt{x^2+y^2}}dz\right]dx\,dy\\
&=\int_0^{2\pi}\cos^2\theta\,d\theta\int_0^1r^4(r-r^2)\,dr
=\frac1{42}\int_0^{2\pi}\cos^2\theta\,d\theta=\frac\pi{42}.
\end{aligned}
\] 【OCR核正】径向积分的上下限和幂次按题面及 Jacobian 恢复。

\subsection{原第4页：证明题1---2}\label{ux539fux7b2c4ux9875ux8bc1ux660eux989812}

\subsubsection{三1. Dirichlet
判别与绝对值比较}\label{ux4e091.-dirichlet-ux5224ux522bux4e0eux7eddux5bf9ux503cux6bd4ux8f83}

记 \(a_n(x)=(-1)^n\)、\(b_n(x)=(x^2+n)/n^2\)。有 \[
\left|\sum_{k=1}^na_k(x)\right|\le1.
\] \(b_n\) 对 \(n\) 单调递减。设 \(M=\max\{|a|,|b|\}\)，则 \[
0\le\sup_{x\in[a,b]}b_n(x)\le\frac{M^2+n}{n^2}\to0.
\] 所以由 Dirichlet 判别法，\(\sum(-1)^n(x^2+n)/n^2\) 在任意有限区间
\([a,b]\) 一致收敛。对每个实数点， \[
\left|(-1)^n\frac{x^2+n}{n^2}\right|\ge\frac1n,
\] 故绝对值级数发散。 【OCR核正】部分和上限为有限的 \(n\)，\(b_n\)
分子为 \(x^2+n\)，不能使用误识的无穷求和或 \(x^n\)。

\subsubsection{三2. 两次使用 Rolle
定理}\label{ux4e092.-ux4e24ux6b21ux4f7fux7528-rolle-ux5b9aux7406}

设 \(f\) 是 \([-1,1]\) 上二阶可导奇函数，\(f(1)=1\)。

（1）令 \(F(x)=f(x)-x\)。由 \(f(0)=0\)，有 \(F(0)=F(1)=0\)。Rolle
定理给出 \(\xi\in(0,1)\)，使 \(F'(\xi)=0\)，即 \(f'(\xi)=1\)。

（2）令 \(G(x)=e^x(f'(x)-1)\)。奇函数的导数为偶函数，所以
\(f'(-\xi)=f'(\xi)=1\)，从而 \(G(-\xi)=G(\xi)=0\)。再用 Rolle 定理，存在
\(\eta\in(-\xi,\xi)\)，使 \[
0=G'(\eta)=e^\eta[f''(\eta)+f'(\eta)-1].
\] 即 \(f''(\eta)+f'(\eta)=1\)。 【OCR核正】将误识为同一个 \(y\)
的不同中值点分别记为 \(\xi,\eta\)，恢复区间关系。

\subsection{原第5页：证明题3---4，综合题1}\label{ux539fux7b2c5ux9875ux8bc1ux660eux989834ux7efcux5408ux98981}

\subsubsection{三3.
换元后对正方形作对称分割}\label{ux4e093.-ux6362ux5143ux540eux5bf9ux6b63ux65b9ux5f62ux4f5cux5bf9ux79f0ux5206ux5272}

\(D=\{x,y\ge0:x+y\le1\}\)。先令 \(t=1-y\)，依原稿换元得 \[
\begin{aligned}
I&=\iint_Df(1-y)f(x)\,dx\,dy\\
&=\int_0^1f(x)\left[\int_0^{1-x}f(1-y)\,dy\right]dx
=\int_0^1f(x)\left[\int_x^1f(t)\,dt\right]dx.
\end{aligned}
\] 换元后的上三角形 \(0\le x\le t\le1\) 与下三角形 \(0\le t\le x\le1\)
关于 \(t=x\) 对称，且 \(f(x)f(t)\) 对交换变量不变，故 \[
\begin{aligned}
2I&=\int_0^1f(x)\left[\int_x^1f(t)\,dt\right]dx
+\int_0^1f(x)\left[\int_0^xf(t)\,dt\right]dx\\
&=\int_0^1f(x)\,dx\int_0^1f(t)\,dt
=\left[\int_0^1f(x)\,dx\right]^2.
\end{aligned}
\] 【OCR核正】换元后的下限为
\(x\)。原图的区域关系以明确不等式保留，未复刻原图。

\subsubsection{三4. Laplace
乘积公式}\label{ux4e094.-laplace-ux4e58ux79efux516cux5f0f}

在所需二阶偏导存在且运算成立时，逐项使用乘积法则： \[
\begin{aligned}
\Delta(uv)
&=\partial_{xx}(uv)+\partial_{yy}(uv)+\partial_{zz}(uv)\\
&=u(v_{xx}+v_{yy}+v_{zz})+v(u_{xx}+u_{yy}+u_{zz})
+2(u_xv_x+u_yv_y+u_zv_z)\\
&=u\Delta v+v\Delta u+2\nabla u\cdot\nabla v.
\end{aligned}
\] 【OCR核正】\(u,v,V\) 混淆的项按逐项求导恢复，梯度间恢复点乘符号。

\subsubsection{四1. 保留 Lagrange
乘数法并核定最小值}\label{ux56db1.-ux4fddux7559-lagrange-ux4e58ux6570ux6cd5ux5e76ux6838ux5b9aux6700ux5c0fux503c}

约束为 \(g_1=x+2y-1=0\)、\(g_2=x^2+2y^2+z^2-1=0\)，求 \(d=x^2+y^2+z^2\)
最小值点。令 \[
\mathcal L=d+\lambda g_1+\mu g_2.
\] 驻点方程为 \[
\begin{cases}
2x+\lambda+2\mu x=0,\\
2y+2\lambda+4\mu y=0,\\
2z+2\mu z=0,\\
x+2y=1,\\
x^2+2y^2+z^2=1.
\end{cases}
\] 若 \(z\ne0\)，则 \(\mu=-1\)，前两式给出 \(\lambda=0,y=0\)，约束迫使
\(x=1,z=0\)，矛盾。因此 \(z=0\)，代回约束，得候选点 \[
(1,0,0),\qquad(-1/3,2/3,0).
\]
约束集闭且有界，两约束梯度在交线上线性无关，因此最值必在这些候选点取得。距离平方分别为
\(1\)、\(5/9\)，故最近点为 \((-1/3,2/3,0)\)。

【OCR核正及助手核算补全】原识别将两个乘数混写为
\(\lambda\)，末点横坐标写成
\(-2/3\)。该点不满足约束，不能作为答案；正确点与打印答案及约束一致。候选点比较补足了最小值判定。原稿相应笔迹仍可参照PDF核对。

\subsection{原第6页：综合题2}\label{ux539fux7b2c6ux9875ux7efcux5408ux98982}

设 \[
f(x)=\int_1^\infty\frac{\cos(xt)}{t(1+t^2)}\,dt,\qquad x\in\mathbb R.
\]

\subsubsection{四2（1）积分关于全实轴参数一致收敛}\label{ux56db21ux79efux5206ux5173ux4e8eux5168ux5b9eux8f74ux53c2ux6570ux4e00ux81f4ux6536ux655b}

【原推导核查】本页尝试用 Dirichlet 判别法，将 \(\cos(xt)\)
的积分界写成不依赖 \(x\) 的 \(2\)。但 \(x\ne0\) 时通常只有 \[
\left|\int_A^B\cos(xt)\,dt\right|\le2/|x|,
\] 不能在 \(x\to0\) 时统一控制；\(x=0\) 时积分为
\(B-A\)。该原思路不足以证明全实轴上一致收敛。

【助手补充修正】使用可积控制函数： \[
\left|\frac{\cos(xt)}{t(1+t^2)}\right|\le\frac1{t(1+t^2)},\qquad t\ge1.
\] 右端在 \([1,\infty)\) 可积，尾部积分趋于零且不依赖
\(x\)，所以原积分在全实轴关于 \(x\) 一致收敛。

\subsubsection{四2（2）用余弦差公式证明一致连续}\label{ux56db22ux7528ux4f59ux5f26ux5deeux516cux5f0fux8bc1ux660eux4e00ux81f4ux8fdeux7eed}

保留手写版通过三角恒等式直接估计的独有推导。对任意
\(x_1,x_2\in\mathbb R\)， \[
\begin{aligned}
|f(x_1)-f(x_2)|
&\le\int_1^\infty
\frac{2|\sin((x_1+x_2)t/2)\sin((x_1-x_2)t/2)|}{t(1+t^2)}\,dt\\
&\le|x_1-x_2|\int_1^\infty\frac{dt}{1+t^2}
=\frac\pi4|x_1-x_2|.
\end{aligned}
\] 给定 \(\varepsilon>0\)，取
\(\delta=4\varepsilon/\pi\)，即可得到一致连续定义要求。

\subsection{局部核对范围}\label{ux5c40ux90e8ux6838ux5bf9ux8303ux56f4}

六页题号、跨页关系及可靠推导已按逐页OCR整理。原图没有复刻。第5页末点横坐标、第6页
Dirichlet
估计需特别参照原笔迹或原论证；正确数值与补充证明已明确归属于助手整理。填空第1题继续保留打印版函数值与手写版导数的题面差异。
\end{document}
